\appendix
\section{Review protocol and reliability rubric}\label{app:protocol}
% =====================================================================
\paragraph{Identification.} The review follows PRISMA~2020. The primary bibliographic source is the
\textbf{OpenAlex} API; we cross-check with the \textbf{Semantic Scholar} Graph API and a structured
\textbf{web ``deep-research''} arm. The review used a
25-query HCC/HCI/qualitative concept-block protocol (the same query strings were issued against every source so the
arms are comparable). Concept blocks combine an LLM term (``LLM'', ``large language model'', ``GPT-4'', ``ChatGPT'',
``Gemini''), an evaluation/agreement term (``LLM-as-a-judge'', ``automatic evaluation'', ``annotation'', ``inter-annotator
agreement'', ``human evaluation''), and an HCC method term (``qualitative coding'', ``thematic analysis'',
``content analysis'', ``crowdsourcing'', ``participant simulation'', ``persona'', ``usability evaluation'', ``design
critique''). Records were deduplicated by DOI and normalized title.

\paragraph{Screening and eligibility.} \emph{Inclusion:} peer-reviewed or preprint studies that empirically compare
LLM judgments to human judgments on the same items and report (or allow computing) an agreement statistic.
\emph{Exclusion:} non-empirical position papers, studies with no human baseline, studies that use LLMs only to
generate (not judge) content, and off-topic retrievals. Title/abstract screening was followed by full-text
eligibility assessment; per-stage counts appear in the PRISMA flow figure (Figure~\ref{fig:s2prisma}).

\paragraph{Reliability verdict (five-level rubric).} Each study is coded against the \emph{human-human agreement
ceiling} for its task:
\begin{itemize}[leftmargin=1.4em,itemsep=1pt]
\item \textbf{1 Validated:} LLM-human agreement is at or above the human-human ceiling (substitution is empirically
justified).
\item \textbf{2 Promising:} agreement approaches the ceiling under specific, reported conditions (model, prompt,
aggregation); conditional substitution.
\item \textbf{3 Mixed:} results are task- or setup-dependent and inconsistent across studies.
\item \textbf{4 Caution:} agreement is consistently and materially below the ceiling.
\item \textbf{5 Unreliable:} agreement is near chance or dominated by bias (position, verbosity, self-preference).
\end{itemize}
We report a reliability score $=6-\text{verdict}$ and synthesize within metric families (e.g.\ Cohen's $\kappa$,
Krippendorff's $\alpha$, Spearman $\rho$, accuracy) rather than pooling incommensurable effect sizes. The full coded corpus (Appendix~\ref{app:s2}) is released for re-coding.

% =====================================================================
\section{Study 3: datasets, judging prompts, and metrics}\label{app:s3detail}
% =====================================================================
\paragraph{Datasets.} Study 2 uses \SOCndat{} subjective social-computing datasets (Table~\ref{tab:cardssoc}); each
is sampled to $n{=}500$ items with a clear majority label (ties dropped), using the per-item annotator distributions
to compute the human ceiling. MultiPICo (2024) is a recent corpus that serves as a post-cutoff condition.

\begin{table}[h]\centering\small
\caption{Study 2 dataset cards. Ceiling $=$ random-annotator-vs-majority accuracy.}
\label{tab:cardssoc}
\begin{tabular}{@{}l l l c c@{}}
\toprule
Dataset & Task & Source & Annot./item & Ceiling \\
\midrule
GoEmotions & sentiment (4-way, authors' grouping) & Reddit (Google) & ${\sim}3$-$5$ & 0.78 \\
SBIC & offensiveness (binary) & Social Bias Frames & ${\sim}3$ & 0.91 \\
HateXplain & hate speech (3-way) & Twitter/Gab & $3$ & 0.83 \\
MultiPICo & irony (binary) & Reddit/Twitter (2024) & $3$-$8$ & 0.82 \\
\bottomrule
\end{tabular}

\end{table}

\paragraph{Judging prompts.} Each model labels every item as a virtual annotator, returning a JSON label on the final
line (parsed with a surface-form-to-code map):
\begin{quote}\small\ttfamily
Emotion: Comment: \{text\} / Overall sentiment: POSITIVE, NEGATIVE, AMBIGUOUS, or NEUTRAL. -> \{"label": ...\}\\[2pt]
Offensive: Post: \{text\} / Is this post offensive (including potentially offensive)? -> \{"label": "offensive"|"not offensive"\}\\[2pt]
Hate: Post: \{text\} / HATE SPEECH (identity-targeted), OFFENSIVE (not group-targeted), or NORMAL. -> \{"label": ...\}\\[2pt]
Irony: \{post+reply\} / Is the REPLY ironic/sarcastic, or not ironic? -> \{"label": "ironic"|"not ironic"\}
\end{quote}

\paragraph{Metrics.} Let an item have $A$ human annotations with majority label $y^\star$. \textbf{Human ceiling} is
the random-annotator-vs-majority accuracy averaged over items, with an annotator-bootstrap 95\% CI. \textbf{Accuracy},
\textbf{Cohen's $\kappa$}, and \textbf{macro-F1} are against $y^\star$; the \textbf{jury} is the per-item panel
majority \citep{verga2024poll}; \textbf{95\% CIs} are nonparametric bootstrap (1000 item resamples). A
\textbf{contamination screen} (one-third-prefix verbatim reproduction) is reported as a conservative check; a
clean-control (Appendix~\ref{app:robust}) shows it false-positives on high-recall models, so we rely on the recent
MultiPICo post-cutoff condition. The run recorded \emph{zero} API errors and a 100\% answer-parse rate across all
\SOCnjudg{} judgments.

\paragraph{Inferential model.} We model each judgment's correctness with an item-level Bayesian logistic GLMM with
random intercepts for item and model, reporting variance components on a common logit scale (item
$\mathrm{SD}_{\text{logit}}=\SOCitemSDlogit$, model $\SOCmodelSDlogit$). With four datasets we do not fit a
dataset-level model. We bin items by human entropy and report the judge-vs-ceiling \emph{gap} per bin to show the
shortfall is not a label-noise artifact (the gap is largest where humans most agree, see the main text).

% =====================================================================
\section{Inter-coder reliability}\label{app:irr}
% =====================================================================
Because the reliability verdict and the role/stance codes embed judgement, we re-coded a held-out random sample of
each corpus with two \emph{independent} coders who saw only the per-study summaries and the rubric (no original
codes, no other context). For the review ($n{=}\IRRstwoN$, a 30\% sample) the reliability verdict reproduced at
Krippendorff $\alpha=\IRRstwoV$, the augment/replace/simulate role at $\alpha=\IRRstwoRole$, and the epistemic stance
at $\alpha=\IRRstwoStance$ (the least reliable code). \emph{Caveat:} the re-coders are
LLM-assisted, like the primary pipeline, so these figures measure rubric \emph{reproducibility} under independent
application, not human inter-rater agreement; the moderate stance $\alpha$ is why we treat stance-based claims as
indicative.

% =====================================================================
\section{Robustness and validity checks}\label{app:robust}
% =====================================================================
\textbf{Prompt robustness.} We re-ran a robustness subset (three models, $n{=}120$ per cell) under
three semantically equivalent prompt paraphrases per task. \emph{Accuracy} is fairly prompt-stable (mean SD across
paraphrases $=\PRmeanSD$), but the discrete five-level \emph{verdict} is exactly stable in only \PRstable{}/\PRtotal{}
cells: most instability is a one-level flip for cells sitting near a band boundary, and the smaller model is the most
sensitive (worst-case acc SD $\PRmaxSD$), the larger models the least. We therefore report verdicts as robust to
\emph{within one level} rather than exact, and flag small-model verdicts on subjective tasks as prompt-sensitive.
\textbf{Contamination clean control.} We ran the verbatim-recall probe on freshly authored, post-cutoff sentences
that no model can have seen. The result is a cautionary one: the high-recall model still scores $\CCmaxOverlap$ mean
overlap (\CCflag) on this guaranteed-clean text, i.e.\ it confabulates plausible continuations regardless of
memorization, while the low-recall model scores LOW. The probe thus has a high false-positive rate for high-recall
models and \emph{cannot} be read as a reliable leakage test; the uniform MEDIUM tiers in Table~\ref{tab:socres} are
largely an artifact of generative fluency. We retain the probe only as a conservative, low-recall-model signal and
base our substantive claims on the disagreement analysis instead.
\textbf{Annotator-subgroup check.} On MultiPICo (which carries annotator demographics), the panel's agreement with
the female- vs male-annotator majority differs by only \MPCgenderGap, i.e.\ we did not detect a large majority-group
bias in this corpus, though demographic metadata is sparse and this check is limited.
\textbf{Ceiling vs annotator count.} The random-annotator-vs-majority ceiling depends on how many annotators an item
has: subsampling ChaosNLI-SNLI (${\sim}100$ annotators) to $A{=}3$ raises its estimated ceiling from
\RTwoablCeilFull{} to \RTwoablCeilThree{} (fewer annotators yield more unanimous majorities). We therefore compute
each dataset's ceiling at its \emph{native} annotator count and caution that ceilings are not directly comparable
across datasets with very different $A$; the disagreement-aware analyses (entropy, $\kappa$) do not depend on this
estimator.

% =====================================================================
\section{Study 2 corpus}\label{app:s2}
% =====================================================================
\footnotesize
\begin{longtable}{@{}l l l l c l@{}}
\caption{Study 2 corpus (55 studies): HCC/HCI/qualitative methods using LLMs vs human input.}\label{tab:study2full}\\
\toprule
Study & Method & Role & Community & Agree & Verdict \\
\midrule
\endfirsthead
\toprule Study & Method & Role & Community & Agree & Verdict \\ \midrule \endhead
Chew et al. (2023) & Content analysis & augment & CSS & -- & Promising \\
Leas et al. (2023) & Content analysis & augment & Health & -- & Mixed \\
Long et al. (2024) & Content analysis & augment & Education & -- & Mixed \\
Dunivin (2025) & Content analysis & augment & CSS & -- & Promising \\
Anderson et al. (2024) & Creativity/ideation & augment & HCI & -- & Caution \\
Gilardi et al. (2023) & Crowdsourced annotation & replace & CSS & -- & Validated \\
Tornberg (2023) & Crowdsourced annotation & replace & CSS & -- & Validated \\
Kuzman et al. (2023) & Crowdsourced annotation & replace & NLP & -- & Mixed \\
Li et al. (2023) & Crowdsourced annotation & augment & NLP & -- & Promising \\
Ostyakova et al. (2023) & Crowdsourced annotation & replace & NLP & -- & Mixed \\
Tabone \& de Winter (2023) & HCI research methods & augment & HCI & -- & Mixed \\
Aubin Le Quere et al. (2024) & HCI research methods & augment & HCI & -- & Mixed \\
Salminen et al. (2024) & Persona generation & augment & HCI & -- & Mixed \\
Jung et al. (2025) & Persona generation & augment & HCI & -- & Mixed \\
Haxvig et al. (2025) & Persona generation & replace & HCI & -- & Caution \\
Zhang et al. (2023) & Qual analysis (general) & augment & Qual-methods & -- & Mixed \\
Wachinger et al. (2024) & Qual analysis (general) & augment & Health & -- & Mixed \\
Li et al. (2024) & Qual analysis (general) & augment & Health & 0.40 & Mixed \\
Mehta et al. (2025) & Qual analysis (general) & augment & Qual-methods & -- & Mixed \\
Davison et al. (2024) & Qual analysis (meta) & augment & CSS & -- & Caution \\
Schroeder et al. (2025) & Qual analysis (meta) & augment & HCI & -- & Mixed \\
Zhang et al. (2024) & Qual analysis (tool) & augment & HCI & -- & Mixed \\
Meng et al. (2024) & Qual analysis (tool) & augment & HCI & -- & Mixed \\
Xiao et al. (2023) & Qual coding (deductive) & augment & HCI & -- & Promising \\
Savelka et al. (2023) & Qual coding (deductive) & augment & NLP & -- & Mixed \\
Zhang et al. (2025) & Qual coding (deductive) & augment & Qual-methods & 0.46 & Mixed \\
authors (2025) & Qual coding (deductive) & augment & Qual-methods & 0.46 & Mixed \\
Gao et al. (2024) & Qual coding (inductive) & augment & HCI & -- & Promising \\
Dunivin (2024) & Qual coding (inductive) & replace & Qual-methods & -- & Promising \\
Ngo et al. (2026) & Qual coding (inductive) & augment & HCI & -- & Mixed \\
Wu et al. (2023) & Social-computing annotation & replace & CSCW & -- & Mixed \\
Zhu et al. (2023) & Social-computing annotation & replace & CSCW & -- & Mixed \\
Rajashekar et al. (2024) & Social-computing annotation & augment & HCI & -- & Mixed \\
Argyle et al. (2023) & Synthetic participants & simulate & CSS & -- & Mixed \\
Bisbee et al. (2023) & Synthetic participants & replace & CSS & -- & Caution \\
Cheng et al. (2023) & Synthetic participants & simulate & NLP & -- & Caution \\
Wang et al. (2024) & Synthetic participants & replace & HCI & -- & Caution \\
Bail (2024) & Synthetic participants & simulate & CSS & -- & Mixed \\
Abbasiantaeb et al. (2024) & Synthetic participants & simulate & IR & -- & Mixed \\
Huang et al. (2025) & Synthetic participants & simulate & CSS & -- & Mixed \\
Salminen et al. (2025) & Synthetic participants & replace & HCI & -- & Caution \\
De Paoli (2023) & Thematic analysis & augment & CSS & -- & Mixed \\
Dai et al. (2023) & Thematic analysis & augment & NLP & -- & Promising \\
Prescott et al. (2024) & Thematic analysis & augment & Health & -- & Mixed \\
Nguyen-Trung (2025) & Thematic analysis & augment & Qual-methods & -- & Mixed \\
Jain et al. (2025) & Thematic analysis & replace & Qual-methods & 0.91 & Promising \\
Castellanos et al. (2025) & Thematic analysis & augment & Health & 0.80 & Promising \\
Sharma et al. (2025) & Thematic analysis & augment & HCI & -- & Promising \\
Wen et al. (2025) & Thematic analysis & augment & Qual-methods & -- & Mixed \\
Kocaballi (2023) & UX/design critique & augment & HCI & -- & Mixed \\
Duan et al. (2024) & UX/design critique & augment & HCI & -- & Mixed \\
Duan et al. (2024) & UX/design critique & augment & HCI & -- & Mixed \\
Duan et al. (2024) & UX/design critique & augment & HCI & -- & Mixed \\
Guerino et al. (2025) & Usability/heuristic eval & replace & HCI & 0.21 & Caution \\
Zhong et al. (2025) & Usability/heuristic eval & replace & HCI & 0.75 & Mixed \\
\bottomrule
\end{longtable}
